\subsection{Compatibility of the dodecaphase and analytic representations of alpha}
\label{subsec:alpha-limite-dos-vias}

Compatibility is established at every depth. The reversible
dodecaphase representation and the analytic representation of the alpha
coordinate
\(E_{21/22}\) are two internal representations with distinct codomains,
produced from the same APP--TRIT--TPK enriched state. The jet
\(E_{21/22}^{(9)}\) is an exact finite witness, but does not by itself
determine the tail. The completion constructed below does not follow from
the bare jet: it receives the shared coordinate of the state and expresses
its analytic image through an explicit recurrence. It therefore does not
constitute an independent causal selection; it constitutes a typed analytic
representation of the same generated point.

This distinction preserves the causal order
\[
 \mathrm{APP}\longrightarrow\mathrm{TRIT}\longrightarrow\mathrm{TPK}
 \longrightarrow\mathcal X_{\alpha}^{\mathrm{enr}}
 \longrightarrow (P,E,\Phi,U_{12}^{\mathrm{sgn}},K)
 \longrightarrow\alpha_{\mathrm{HMT}}.
\]
The coordinates \(P,E,\Phi\) and the record \(K\) are prior outputs of the
same enriched state. Their use in the closure of alpha is therefore a
subsequent internal composition, not an injection of conventional
constants.

For \(r\in6\mathbb N\), \(r\ge36\), let
\(\widehat{\mathcal X}^{\mathrm{enr},\alpha}_r\) be the sector of truncated
states extending the internal alpha seed fixed at \(R_{36}\). Write
\[
 x_r^\alpha=(\widetilde x_r,\chi_\alpha,
              \varepsilon_r^\alpha,N_r^\alpha),
\]
where \(\chi_\alpha\) is the internal character of the sector, not the
expressed scalar value. Each state preserves residue, quotient, carry, sheet,
orientation, phase, route, cylinder, boundary, signature, record, memory,
survival, provenance, and depth. If
\[
 d_r^\alpha=\operatorname{Dig}_{729}(\varepsilon_r^\alpha),\qquad
 \varepsilon_{r+6}^\alpha=729\varepsilon_r^\alpha-d_r^\alpha,
 \qquad N_{r+6}^\alpha=729N_r^\alpha+d_r^\alpha,
\]
the effective transition is
\[
 \mathcal U_r^\alpha=
 \operatorname{Upd}_r^\alpha\circ
 \operatorname{Tra}_r^\alpha\circ
 \operatorname{Sel}_r^\alpha,
 \qquad
 \operatorname{Ext}^{\alpha,\to}_{r+6,r}
 =\operatorname{graph}(\mathcal U_r^\alpha).
\]
The functional character of these extensions, with the seed \(R_{36}\) fixed,
produces a unique compatible section: phase returns while memory and
depth advance. The global domain of both determinations is
\begin{equation}
 \mathcal X_\alpha^{\mathrm{enr}}:=\widehat\Omega_\alpha
 :=\varprojlim_{r\ge36}
 \bigl(\widehat{\mathcal X}^{\mathrm{enr},\alpha}_r,
       \operatorname{Ext}^{\alpha,\to}_{r+6,r}\bigr).
 \label{eq:alpha-dominio-enriquecido-completo}
\end{equation}

\subsubsection{Dodecaphase closure as a complete section}

For \(N\geq1\), set
\[
 \mathcal W_N:=\{0,\ldots,999\}^{N},
 \qquad
 \mathsf T_{\alpha,N}(\mathsf s):=(A_1,\ldots,A_N),
 \qquad \mathsf s\in\mathcal X_{\alpha}^{\mathrm{enr}},
\]
where the blocks \(A_j\) are obtained through
\eqref{eq:alpha-recurrencia} with the compatible boundary
\(c_{N+1}^{\infty}=\lfloor R_{N+1}(Z)\rfloor\). For \(N'\geq N\), let
\(\rho_{N',N}:\mathcal W_{N'}\to\mathcal W_N\) be restriction to the
first \(N\) blocks.

\begin{theorem}[Complete dodecaphase output and projective compatibility]
\label{thm:alpha-salida-dodecafasica-completa}
The collection \((\mathsf T_{\alpha,N})_{N\geq1}\) is compatible:
\begin{equation}
 \rho_{N',N}\circ\mathsf T_{\alpha,N'}
 =\mathsf T_{\alpha,N}
 \qquad(N'\geq N).
 \label{eq:alpha-via-a-proyectiva-exacta}
\end{equation}
Its canonical cylinders
\begin{equation}
 C_N^A(\mathsf s):=
 \left[
  \sum_{j=1}^{N}A_j1000^{-j},
  \sum_{j=1}^{N}A_j1000^{-j}+1000^{-N}
 \right)\cap I_\alpha
 \label{eq:alpha-cilindros-via-a-exactos}
\end{equation}
are nested, have diameter at most \(1000^{-N}\), and satisfy
\begin{equation}
 \bigcap_{N\geq1}C_N^A(\mathsf s)
 =\{\alpha_A(\mathsf s)\},
 \qquad
 \alpha_A
 =\pi_{\mathrm{HMT}}+e_{\mathrm{HMT}}-\varphi_{\mathrm{HMT}}
  -4-\kappa_{\mathrm{per}}.
 \label{eq:alpha-salida-a-completa}
\end{equation}
Consequently, the first route expresses a unique infinite section, denoted
by \(\alpha_{\mathrm{HMT}}:=\alpha_A\).
\end{theorem}

\begin{proof}
Proposition~\ref{prop:alpha-normalizacion} provides the unique
canonical expansion and the unique bounded carry sequence. The local identity
\eqref{eq:alpha-identidad-local} and its telescoping form
\eqref{eq:alpha-telescopia} preserve the remote boundary; Proposition
\ref{ampl-alfa-prop-frontera} proves that restricting an extended execution
recovers precisely the previous prefix. This gives
\eqref{eq:alpha-via-a-proyectiva-exacta}. Membership of the complete sum
in every cylinder follows from the tail identity, and
\(\operatorname{diam}C_N^A\leq1000^{-N}\to0\); the intersection contains
a single point. Finally, \eqref{eq:alpha-identidad-completa} identifies this
point with the expression in \eqref{eq:alpha-salida-a-completa}. All its
coefficients and signs are fixed before alpha is expressed.
\end{proof}

The window \(\alpha_{12}^{(0)}\), closed with \(c_{13}=0\), remains an
exact finite calculation; it must not be confused with the restriction of
the complete section, whose compatible boundary satisfies \(c_{13}^{\infty}=-1\).
In particular,
\begin{equation}
 \rho_{36,12}(\alpha_{A,36})=\alpha_{A,12}
 \label{eq:alpha-rho-36-12-proyectiva}
\end{equation}
is an equality of compatible words, not an equality between words of
different lengths or an identification of the isolated window with the
limit.

\subsubsection{Exact analytic characterization of order nine}

Let
\[
 \mathbb K:=\mathbb Q\!\left(\pi_{\mathrm{HMT}},
                  \log_{10}\frac{10}{9}\right),
 \qquad
 E_9(x):=E_{21/22}^{(9)}(x)=P_9(x)-\delta .
\]
Equations \eqref{eq:alpha-firma-coeficientes} and
\eqref{eq:alpha-t79} construct \(P_9=P_6+T_{7:9}\) from the state's
signature and the nilpotent resolvent;
Proposition~\ref{prop:alpha-raiz-jet} proves that \(E_9\) has a unique
simple root \(\xi_9\in I_\alpha\), isolated by the rational interval
\eqref{eq:alpha-intervalo-jet}.

Finite bidirectionality can be expressed without introducing an
analytic tail. Define
\begin{align}
 B_9(x):={}&\delta+\frac74x^2-\frac{x^4}{20}
 +\frac{2x^5}{21}+\frac{x^6}{46}+\frac{x^7}{120}
 -\frac{x^8}{45}-\frac{2x^9}{495},
 \label{eq:alpha-b9-finito}\\
 q_x(T):={}&xT^2-B_9(x)T+x^3,\qquad x\in I_\alpha=(0,10^{-2}),
 \label{eq:alpha-q-involucion-finita}\\
 J_x:{\mathbb R}^{\times}\longrightarrow{\mathbb R}^{\times},
 \qquad J_x(T):={}&\frac{x^2}{T}.
\end{align}

\begin{proposition}[Finite quadratic equivalence and involution]
\label{prop:alpha-involucion-finita-exacta}
For every \(x\in I_\alpha\) and every \(T\in\mathbb R^\times\),
\begin{align}
 \frac{q_x(2\pi_{\mathrm{HMT}})}{2\pi_{\mathrm{HMT}}}
 &=P_9(x)-\delta=E_9(x),
 \label{eq:alpha-q-equivale-p9}\\
 J_x(J_x(T))&=T,
 \qquad
 T^2q_x(J_x(T))=x^2q_x(T).
 \label{eq:alpha-j-involucion-exacta}
\end{align}
Let \(Z_x=\{T\in\mathbb R^\times:q_x(T)=0\}\). Consequently,
\(J_x|_{Z_x}:Z_x\to Z_x\) is involutive. When \(E_9(x)=0\),
it exchanges the roots \(2\pi_{\mathrm{HMT}}\) and
\(x^2/(2\pi_{\mathrm{HMT}})\). This is an exact involution on the scalar
fibre of the order-nine jet; it does not act on the enriched state.
\end{proposition}

\begin{proof}
Dividing \(q_x(T)\) by \(T\) and substituting
\(T=2\pi_{\mathrm{HMT}}\) reproduces the definition of
\(P_9(x)-\delta\) term by term. The first identity in
\eqref{eq:alpha-j-involucion-exacta} is immediate. For the second,
\[
 T^2q_x(x^2/T)
 =x^5-B_9(x)x^2T+x^3T^2=x^2q_x(T).
\]
Thus \(J_x\) preserves the root set and, by Viète's formula,
exchanges the two roots whose product is \(x^2\).
\end{proof}

The finite decimal witness thus retains its exact force:
\begin{equation}
 \operatorname{Tr}_9(\xi_9^{-1})
 =\operatorname{Tr}_9\!\left((\alpha_{12}^{(0)})^{-1}\right)
 =137.035999084.
 \label{eq:alpha-testigo-finito-nueve}
\end{equation}
Equality \eqref{eq:alpha-testigo-finito-nueve} certifies the finite
characterization and the preceding involution. It does not turn the root
of a truncated polynomial into the root of every compatible extension.

\subsubsection{The jet alone does not determine an analytic extension}

Let
\[
 j_9:\mathbb K[[x]]\longrightarrow \mathbb K[x]/(x^{10})
\]
be truncation, and for \(h\in\mathbb K[x]\) set
\begin{equation}
 E_h(x):=E_9(x)+x^{10}h(x).
 \label{eq:alpha-familia-prolongaciones-mismo-jet}
\end{equation}
All these polynomials satisfy \(j_9(E_h)=j_9(E_9)\).

\begin{theorem}[Non-uniqueness of the root from the jet]
\label{thm:alpha-no-unicidad-prolongacion-jet}
There is no map
\[
 \mathfrak r_9:\mathbb K[x]/(x^{10})\longrightarrow I_\alpha
\]
that assigns, from the jet alone, the positive root of every
polynomial extension of \(E_9\) having a unique root in
\(I_\alpha\). In particular, \(E_0\) and \(E_{1/729}\) have the same jet
and distinct positive roots.
\end{theorem}

\begin{proof}
The bound established in the proof of
\ref{prop:alpha-raiz-jet} gives \(E_9'(x)>6.24\) on
\([0,10^{-2}]\). For \(t\geq0\), write
\(E_t(x)=E_9(x)+tx^{10}\). Then
\[
 E_t'(x)=E_9'(x)+10tx^9>6.24,
\]
whereas \(E_t(0)=-\delta<0\) and
\(E_t(10^{-2})>E_9(10^{-2})>0\). Each \(E_t\) therefore has a unique
positive root \(\eta(t)\in I_\alpha\). For \(t=1/729\),
\[
 E_{1/729}(\xi_9)=\frac{\xi_9^{10}}{729}>0;
\]
since \(E_{1/729}\) is strictly increasing and changes sign only once,
its root satisfies \(\eta(1/729)<\xi_9=\eta(0)\). However,
\(j_9(E_{1/729})=j_9(E_0)\). A map that factored the root through
\(j_9\) would have to assign the same point to both polynomials, a contradiction.
\end{proof}

The variation is not merely existential. The implicit function theorem
applies because \(\partial_xE_t(\eta(t))>6.24\), and gives
\begin{equation}
 \eta'(t)=-
 \frac{\eta(t)^{10}}
 {E_9'(\eta(t))+10t\eta(t)^9}<0.
 \label{eq:alpha-variacion-raiz-prolongacion}
\end{equation}
Thus coefficients beyond degree nine are proved to change the root
even while leaving the entire jet constructed through that degree invariant.
This conclusion does not deny the tail produced by the enriched state:
it proves only that the root does not factor through the bare truncation
\(j_9\). Choosing a tail \(h\) from an external conventional root
would be circular. The completion below consults neither that root nor a
metrological table: it directly evaluates the regional and terminal coordinates
of the state. Since that evaluation agrees algebraically with the integer
representation, the result is a correlated representation, not an assertion
of independence between two selectors.

% Correlated analytic representation of alpha: action, convergence, and
% intervals.
\subsubsection{Analytic representation of the alpha coordinate}
\label{subsec:alpha-lector-completo-rev08}

The preceding jet control requires the tail to be exhibited: merely naming a
complete reader is insufficient. Let \(\mathsf s\in\mathcal X_\alpha^{\mathrm{enr}}\). Before
decimal normalization, the same state already expresses the precoordinate
\begin{equation}
 a(\mathsf s):=p(\mathsf s)+e_0(\mathsf s)-f(\mathsf s)
                 -\kappa_{\mathrm{per}}(\mathsf s)\in I_\alpha,
 \label{eq:alpha-precoordenada-comun}
\end{equation}
with \(p=\pi_{\mathrm{HMT}}-3\),
\(e_0=e_{\mathrm{HMT}}-2\) and
\(f=\varphi_{\mathrm{HMT}}-1\). Equivalently,
\(a=\pi_{\mathrm{HMT}}+e_{\mathrm{HMT}}-
\varphi_{\mathrm{HMT}}-4-\kappa_{\mathrm{per}}\). Conventional
\(\alpha\) is not introduced here: \(p,e_0,f\) are the three regional characters already
generated, and \(\kappa_{\mathrm{per}}\) is the rational reading of the terminal
register. Integer closure applies to
\eqref{eq:alpha-precoordenada-comun} the normalization of
\eqref{eq:alpha-recurrencia}; the following chart expresses a compatible analytic
representation on the same precoordinate. Its dependency graph
receives \((p,e_0,f,\kappa_{\mathrm{per}})\) directly: it contains no
edge \(\alpha_A\to\lambda\), although the equality proved below
subsequently permits the two representation names to be identified.

Write \(q=1/729\), let \(F_{\mathsf s}\) be the ordered field generated by
the coordinates of \(\mathsf s\), and put
\(E_{9,\mathsf s}(X)=E_{21/22}^{(9)}(\mathsf s;X)\). Define the linking
coefficient
\begin{equation}
 \lambda(\mathsf s):=
 -\frac{E_{9,\mathsf s}(a(\mathsf s))
         \bigl(1-q\,a(\mathsf s)^3\bigr)}{a(\mathsf s)^{10}}.
 \label{eq:alpha-lambda-estado}
\end{equation}
All its arguments belong to the state before metrological recognition.
In particular, \(\lambda\) is not obtained by fitting a target
decimal string. Equation~\eqref{eq:alpha-lambda-estado} must be read with
its exact type: once the first representation of the state has produced
\(a(\mathsf s)\), it fixes the unique prolongation in the geometric class
defined below. It is not a second independent selection of
\(a(\mathsf s)\), but a correlated analytic representation of the same
enriched state.

Indeed, for \(\mu\in F_{\mathsf s}\), put
\[
 E_{\mathsf s,\mu}(X)
 :=E_{9,\mathsf s}(X)+\mu\frac{X^{10}}{1-qX^3}.
\]
Since \(0<a(\mathsf s)<10^{-2}\) and \(qa(\mathsf s)^3<1\), we have
\begin{equation}
 E_{\mathsf s,\mu}(a(\mathsf s))=0
 \quad\Longleftrightarrow\quad
 \mu=\lambda(\mathsf s).
 \label{eq:alpha-unicidad-lambda-clase-geometrica}
\end{equation}
Thus the coefficient seed contains no freedom once the state and the
class of \(q\)-geometric prolongations have been fixed. This uniqueness is relative to the
declared class; it does not assert that the ninth-order jet alone selects a
tail among all possible analytic series.

The local action on blocks of three coefficients is
\begin{equation}
 \mathcal C_\alpha:F_{\mathsf s}^{3}\longrightarrow F_{\mathsf s}^{3},
 \qquad \mathcal C_\alpha(u,v,w)=q(u,v,w),
 \qquad c^{(0)}=(\lambda(\mathsf s),0,0).
 \label{eq:alpha-accion-coeficientes}
\end{equation}
Consequently, for every \(n\geq0\),
\begin{equation}
 b_{10+3n}=q^n\lambda(\mathsf s),\qquad
 b_{11+3n}=b_{12+3n}=0.
 \label{eq:alpha-recurrencia-coeficientes}
\end{equation}
This is the rule that computes every coefficient beyond degree nine. No
free parameter remains at each prolongation.

For \(M\geq0\), let
\begin{equation}
 E_{\mathsf s,M}(X):=E_{9,\mathsf s}(X)
 +\lambda(\mathsf s)\sum_{n=0}^{M}q^nX^{10+3n}.
 \label{eq:alpha-truncamientos-completos}
\end{equation}
The complete function is
\begin{equation}
 E_{\mathsf s,\infty}(X)
 :=E_{9,\mathsf s}(X)
 +\lambda(\mathsf s)\frac{X^{10}}{1-X^3/729}.
 \label{eq:alpha-funcion-completa-explicita}
\end{equation}

To formulate the projective system without implicit types, fix
\(d_M=12+3M\) and define the jet fibration
\begin{equation}
 \mathfrak J_M:=
 \coprod_{\mathsf s\in\mathcal X_\alpha^{\mathrm{enr}}}
 \{\mathsf s\}\times F_{\mathsf s}[X]/(X^{d_M+1}).
 \label{eq:alpha-espacio-jets-rev08}
\end{equation}
If \(M'\ge M\), restriction is
\begin{equation}
 \tau_{M',M}(\mathsf s,[P]_{d_{M'}})
 :=(\mathsf s,[P]_{d_M}).
 \label{eq:alpha-truncamiento-jets-rev08}
\end{equation}

\begin{proposition}[Naturality of the analytic recurrence]
\label{prop:alpha-naturalidad-recurrencia-explicita}
Let
\[
 \Gamma_M(\mathsf s)
 =\bigl(\mathsf s,[E_{\mathsf s,M}]_{d_M}\bigr).
\]
If \(M'\geq M\), the truncation \(\tau_{M',M}\) satisfies
\begin{equation}
 \tau_{M',M}\circ\Gamma_{M'}=\Gamma_M,
 \qquad
 \Gamma_{E21}(\mathsf s)=\varprojlim_M\Gamma_M(\mathsf s).
 \label{eq:alpha-naturalidad-explicita}
\end{equation}
Projection onto the three new degrees is precisely
\(\mathcal C_\alpha^{M+1}(c^{(0)})\).
\end{proposition}

\begin{proof}
Equation~\eqref{eq:alpha-recurrencia-coeficientes} fixes the block of
degrees \(10+3n,11+3n,12+3n\). Passing from \(M\) to \(M+1\) adds only
the next block; truncating it recovers the previous polynomial literally.
Compatibility for \(M'\geq M\) follows by composition. Thus the
inverse limit is not a section chosen from the bare jet, but the unique
orbit of \(c^{(0)}\) under \(\mathcal C_\alpha\).
\end{proof}

\subsubsection{Uniform convergence, a simple root, and rational intervals}
\label{subsec:alpha-cotas-completas-rev08}

The regional cylinders and periodic register provide bounds through
rational interval arithmetic. For each coordinate
\[
 c\in\{\pi_{\mathrm{HMT}},e_{\mathrm{HMT}},\varphi_{\mathrm{HMT}}\},
\]
let \(P_c\) be its certified thousand-digit nonadic prefix. The datum used
is not the rational \(P_c\) treated as the complete value, but the half-open
cylinder and its outer closure
\[
 \mathcal C_c^{(1000)}=[P_c,P_c+\varepsilon),
 \qquad
 \overline{\mathcal C}_c^{(1000)}=[P_c,P_c+\varepsilon],
 \qquad \varepsilon=10^{-1000}.
\]
Directed arithmetic operates conservatively with these closures. If
\(P_\pi,P_e,P_\varphi\) denote the three prefixes, this gives
\begin{equation}
\begin{aligned}
 a_-&:=P_\pi+P_e-(P_\varphi+\varepsilon)-4-
       \kappa_{\mathrm{per}},\\
 a_+&:=(P_\pi+\varepsilon)+(P_e+\varepsilon)-P_\varphi-4-
       \kappa_{\mathrm{per}},\\
 A_{1000}&:=[a_-,a_+],
 \qquad a(\mathsf s)\in[a_-,a_+]=A_{1000},
 \qquad \operatorname{diam}A_{1000}=3\varepsilon.
\end{aligned}
\label{eq:alpha-propagacion-colas-rev08}
\end{equation}
The natural interval extension of
\eqref{eq:alpha-lambda-estado} subsequently transports \(A_{1000}\), the cylinder
of \(\pi_{\mathrm{HMT}}\), and the directed rational bound on \(\delta\) into a
closed interval \(\Lambda_{1000}\ni\lambda(\mathsf s)\); none of the three
tails is frozen. In particular,
\begin{equation}
\begin{aligned}
 \frac{7297352569283800997285105472380662}{10^{36}}
 &<a(\mathsf s)<
 \frac{7297352569283800997285105472380663}{10^{36}},\\
 -8.711\mathbin{\times}10^{-6}
 &<\lambda(\mathsf s)<-8.709\mathbin{\times}10^{-6}.
\end{aligned}
 \label{eq:alpha-cotas-precoordenada-lambda}
\end{equation}
These inequalities use certified prefixes with their rational tail bound;
they do not receive an alpha table. The package verifier reproduces them
from the three nonadic outputs and the terminal chart.

Let
\[
 r:=\frac{10^{-6}}{729}<1.372\mathbin{\times}10^{-9}.
\]
For \(0\leq x\leq10^{-2}\), the tail beyond \(M\) satisfies
\begin{equation}
 \left|E_{\mathsf s,\infty}(x)-E_{\mathsf s,M}(x)\right|
 \leq
 8.711\mathbin{\times}10^{-26}\,\frac{r^{M+1}}{1-r}.
 \label{eq:alpha-cota-cola-uniforme}
\end{equation}
Moreover,
\begin{equation}
 \left|\frac{\mathrm d}{\mathrm dx}
  \left(\lambda\frac{x^{10}}{1-x^3/729}\right)\right|
 \leq 8.711\mathbin{\times}10^{-6}
 \left(\frac{10\,10^{-18}}{1-r}
 +\frac{3\,10^{-24}/729}{(1-r)^2}\right)
 <8.712\mathbin{\times}10^{-23}.
 \label{eq:alpha-cota-derivada-cola}
\end{equation}
The derivative tail beyond block \(M\) also admits the
depth-dependent estimate
\begin{equation}
 \sup_{0\le x\le10^{-2}}
 \left|E'_{\mathsf s,\infty}(x)-E'_{\mathsf s,M}(x)\right|
 \le 8.711\mathbin{\times}10^{-24}\,r^{M+1}
 \left(\frac{13+3M}{1-r}+\frac{3r}{(1-r)^2}\right)
 \xrightarrow[M\to\infty]{}0.
 \label{eq:alpha-cota-resto-derivadas-rev08}
\end{equation}

\begin{theorem}[Limit function and unique simple root of the correlated chart]
\label{thm:alpha-raiz-completa-explicita}
The sequence \(E_{\mathsf s,M}\) converges uniformly, together with its
derivatives, to \(E_{\mathsf s,\infty}\) on
\(\overline I_\alpha=[0,10^{-2}]\). We have
\begin{equation}
 E_{\mathsf s,\infty}(a(\mathsf s))=0,
 \qquad
 \inf_{x\in\overline I_\alpha}
 E'_{\mathsf s,\infty}(x)>6.239.
 \label{eq:alpha-raiz-y-derivada-completas}
\end{equation}
Consequently, \(a(\mathsf s)\) is the unique simple root of the complete
function in \(I_\alpha\).
\end{theorem}

\begin{proof}
The ratio of two consecutive terms of the tail is bounded by \(r<1\),
which yields \eqref{eq:alpha-cota-cola-uniforme}; termwise differentiation
and summation of the two geometric series yield
\eqref{eq:alpha-cota-derivada-cola} and
\eqref{eq:alpha-cota-resto-derivadas-rev08}. The jet proof establishes
\(E'_{9,\mathsf s}>6.24\) on \(\overline I_\alpha\); hence the
complete perturbation leaves the derivative above \(6.239\).
Finally, substituting \eqref{eq:alpha-lambda-estado} into
\eqref{eq:alpha-funcion-completa-explicita} makes the two terms evaluated at
\(a(\mathsf s)\) cancel exactly. The function is strictly
increasing; its unique root is therefore simple and coincides with this precoordinate.
\end{proof}

To make the identification effective at every depth, directed endpoint
evaluations first certify
\[
 E_{\mathsf s,\infty}(0)<0<
 E_{\mathsf s,\infty}(10^{-2}),
 \qquad
 E'_{\mathsf s,\infty}\geq\mu:=6239/1000>c:=6
 \quad\hbox{on }[0,10^{-2}].
\]
Continuity ensures the existence of a root in the interval, and the second inequality
ensures uniqueness. Only after establishing this existence, monotonicity, and the
invariant that the bracket contains the root is the mean-value-theorem
fallback applied. Start with \(D_0=[0,10^{-2}]\). If
\(D_j=[r_j,s_j]\), let \(m_j=(r_j+s_j)/2\), and let
\([F_j^-,F_j^+]\) be a directed rational enclosure of
\(E_{\mathsf s,\infty}(m_j)\), obtained by propagating the cylinders of
\(\pi_{\mathrm{HMT}},e_{\mathrm{HMT}},\varphi_{\mathrm{HMT}}\),
\(\delta\), and \(\lambda\), with
\begin{equation}
 0\leq F_j^+-F_j^-\leq \frac{c}{4}\,(s_j-r_j).
 \label{eq:alpha-anchura-evaluacion-rescate-rev08}
\end{equation}
When \(F_j^+<0\), retain the right half;
when \(F_j^->0\), retain the left half. If
\(F_j^-\leq0\leq F_j^+\), no attempt is made to decide whether the midpoint
is exactly the root. Let \(w_j=s_j-r_j\). The uniform bound
\(E'_{\mathsf s,\infty}\geq c=6\) and the mean value theorem imply that
any root contained in \(D_j\) belongs to
\begin{equation}
 D_{j+1}=D_j\cap
 \left[m_j-\frac{w_j}{4},
             m_j+\frac{w_j}{4}\right].
 \label{eq:alpha-rescate-derivada-rev08}
\end{equation}
Indeed, let \(\xi_j\) be the root and let
\(y_j=E_{\mathsf s,\infty}(m_j)\in[F_j^-,F_j^+]\). Since the enclosure contains
zero and satisfies \eqref{eq:alpha-anchura-evaluacion-rescate-rev08}, we have
\(|y_j|\leq cw_j/4\). The mean value theorem gives
\(|m_j-\xi_j|\leq |y_j|/c\leq w_j/4\), proving
\eqref{eq:alpha-rescate-derivada-rev08}. The argument includes the case
\(y_j=0\): the root then lies at the midpoint, but the algorithm does not
need to recognize this equality in order to retain it.

If the available enclosure does not satisfy
\eqref{eq:alpha-anchura-evaluacion-rescate-rev08}, the source cylinders are first
refined and evaluation is repeated; no sign is declared and no
insufficient contraction is accepted. In each of the three branches, the new
interval retains the root and has diameter at most \(w_j/2\). For every
finite target depth, directed refinement permits the width bound to be met
without an uncertified comparison of a real number with zero.
Monotonicity proves inductively
\begin{equation}
 a(\mathsf s)\in D_{j+1}\subseteq D_j,\qquad
 \operatorname{diam}D_{j+1}\leq
 \tfrac12\operatorname{diam}D_j,
 \qquad \operatorname{diam}D_j\longrightarrow0,
 \label{eq:alpha-biseccion-racional}
\end{equation}
where the inequality is a contraction bound, not an imposed equality.
The executable certificate applies
this rule at twenty-four base-thousand levels and verifies that the complete outer
closure \(A_{1000}\), not merely its lower endpoint, lies within
every expressed cylinder. Its negative tests reject an excessively wide
enclosure, a disordered enclosure, a point other than the midpoint, a domain without
a sign change, and a degenerate bracket. Coinductive prolongation permits
a thousand to be replaced by any required depth.

Specifically, for \(S_N=1000^N\), the verifier computes
\begin{equation}
 j_N^-:=\lfloor S_Na_-\rfloor,
 \qquad j_N^+:=\lfloor S_Na_+\rfloor
 \label{eq:alpha-indices-cilindro-dirigidos-rev08}
\end{equation}
and expresses the level only if \(j_N^-=j_N^+\). This equality, which includes the
upper endpoint of the outer closure, proves
\(A_{1000}\subset[j_N^-/S_N,(j_N^-+1)/S_N)\) and prevents an endpoint lying
exactly on the right boundary from being assigned to the left cell.

Let \(C_N^A(\mathsf s)\) be the integer cylinder of
\eqref{eq:alpha-cilindros-via-a-exactos}. Recursively choose
\(m(N)>m(N-1)\) such that
\(\operatorname{diam}D_{m(N)}<1000^{-N}/4\), and define
\begin{equation}
I_{B,N}:=D_{m(N)}\cap C_N^A(\mathsf s).
 \label{eq:alpha-intervalos-b-completos}
\end{equation}
Define
\(\ell_N:=\inf I_{B,N}\) and \(u_N:=\sup I_{B,N}\). Both are rational;
the interval is nonempty because both factors contain
\(a(\mathsf s)\). Since the cylinder is half-open, \(u_N\) may not
belong to \(I_{B,N}\). Evaluation at this endpoint is understood through
the continuous extension of \(E_{\mathsf s,\infty}\) to the outer closure. Thus,
\begin{equation}
 I_{B,N+1}\subseteq I_{B,N}\subseteq C_N^A(\mathsf s),
 \qquad \operatorname{diam}I_{B,N}\longrightarrow0,
 \qquad E_{\mathsf s,\infty}(\ell_N)\leq0\leq
 E_{\mathsf s,\infty}(u_N).
 \label{eq:alpha-inclusion-signos-completa}
\end{equation}
Bound~\eqref{eq:alpha-cota-cola-uniforme} permits a truncation to be
chosen that preserves the sign at every endpoint where the complete function
does not vanish. If an endpoint coincides with the root, the identity of
the complete function and the membership already proved are used; the same
sign is not attributed to its truncations. Indeed, the exact remainder is
\[
 E_{\mathsf s,\infty}(x)-E_{\mathsf s,M}(x)
 =\lambda(\mathsf s)x^{10}
   \frac{(qx^3)^{M+1}}{1-qx^3}.
\]
On the positive domain considered, where \(\lambda(\mathsf s)<0\), this
gives at the root
\[
 E_{\mathsf s,M}(a(\mathsf s))
 =-\lambda(\mathsf s)a(\mathsf s)^{10}
   \frac{(qa(\mathsf s)^3)^{M+1}}{1-qa(\mathsf s)^3}>0
 \qquad(M<\infty).
\]
Consequently, certification of the weak signs in
\eqref{eq:alpha-inclusion-signos-completa} refers to
\(E_{\mathsf s,\infty}\). Its evaluation uses the complete function or the
truncation accompanied by the signed remainder. This covers both the included
infimum and the potentially excluded supremum of the cylinder.

\begin{proposition}[Compatibility squares and common cylinders]
\label{prop:alpha-cuadrados-compatibilidad}
The restrictions of words and analytic representations commute:
\[
\begin{array}{ccc}
 \mathcal X_\alpha^{\mathrm{enr}}&
 \xrightarrow{\ \mathsf T_{\alpha,N'}\ }&\mathcal W_{N'}\\[1mm]
 \Big\Vert&&\Big\downarrow\rho_{N',N}\\[1mm]
 \mathcal X_\alpha^{\mathrm{enr}}&
 \xrightarrow{\ \mathsf T_{\alpha,N}\ }&\mathcal W_N
\end{array}
\qquad
\begin{array}{ccc}
 \mathcal X_\alpha^{\mathrm{enr}}&
 \xrightarrow{\ \Gamma_{M'}\ }&\mathfrak J_{M'}\\[1mm]
 \Big\Vert&&\Big\downarrow\tau_{M',M}\\[1mm]
 \mathcal X_\alpha^{\mathrm{enr}}&
 \xrightarrow{\ \Gamma_M\ }&\mathfrak J_M .
\end{array}
\]
The collection \(I_{B,N}\) provides the effective connection between both
squares and concretely satisfies
\(I_{B,N}\subseteq C_N^A\) for every \(N\).
\end{proposition}

\begin{proof}
The first square is the compatibility of integer normalization; the second
is Proposition~\ref{prop:alpha-naturalidad-recurrencia-explicita}. The
inclusion follows from definition~\eqref{eq:alpha-intervalos-b-completos},
and nonemptiness and signs follow from
Theorem~\ref{thm:alpha-raiz-completa-explicita}. No step identifies the
root of a truncation with that of an arbitrary tail.
\end{proof}

Now define, without leaving the prefix reader implicit,
\begin{equation}
 \alpha_B(\mathsf s):=
 \operatorname*{arg\,zero}_{x\in I_\alpha}E_{\mathsf s,\infty}(x),
 \qquad
 \operatorname{pref}_N(x):=
 \frac{\lfloor1000^Nx\rfloor}{1000^N},
 \qquad
 \alpha_{B,N}:=\operatorname{pref}_N(\alpha_B).
 \label{eq:alpha-definicion-b-y-prefijos-rev08}
\end{equation}

\begin{theorem}[Projective equality of the two correlated representations]
\label{prop:alpha-dos-limites}
For every \(N\), the integer route and the analytic route express the same prefix:
\begingroup
\renewcommand{\theequation}{A7\(_0\)}
\begin{equation}
 \boxed{\alpha_{A,N}=\operatorname{pref}_N(\alpha_A)
 =\operatorname{pref}_N(\alpha_B)=\alpha_{B,N}.}
 \label{eq:alpha-dos-vias-cada-nivel}
\end{equation}
\endgroup
In the limit,
\begingroup
\renewcommand{\theequation}{A8\(_0\)}
\begin{equation}
 \boxed{\alpha_A
 =\operatorname*{arg\,zero}_{x\in I_\alpha}
       E_{\mathsf s,\infty}(x)
 =\alpha_B=\alpha_{\mathrm{HMT}}.}
 \label{eq:alpha-dos-vias-limite-coinductivo}
\end{equation}
\endgroup
\end{theorem}

\begin{proof}
Integer normalization expresses exactly \(a(\mathsf s)\).
Theorem~\ref{thm:alpha-raiz-completa-explicita} proves that the analytic
chart represents the same point as its unique simple root. The
half-open cylinders \(C_N^A\) fix a canonical prefix, and
\(I_{B,N}\subseteq C_N^A\) contains \(\alpha_B=\alpha_A\). Their
prefixes therefore coincide for every \(N\); nesting and diameters tending to
zero yield equality of the limits. Metrological
recognition occurs afterwards.
\end{proof}

The analytic function constitutes a correlated representation of the same
state, not a second derivation causally independent of integer closure.
This precision does not diminish the equality proved: it identifies its
exact mechanism, prevents a tail that the jet does not determine from being attributed to it,
and avoids promoting a conventional value to a generator.

\endinput

% The following block is preserved solely as an editorial antecedent of
% REV04. \endinput prevents it from forming part of the active source or the REV08 PDF.
\subsubsection{Complete reader, naturality, and common cylinders}

Let \(F_{\mathsf s}\) be the coefficient field of the signature of
\(\mathsf s\in\mathcal X_\alpha^{\mathrm{enr}}\), and let \(d_m=6+3m\).
Define
\[
 \mathfrak J_m:=\coprod_{\mathsf s\in\mathcal X_\alpha^{\mathrm{enr}}}
 \{\mathsf s\}\times F_{\mathsf s}[X]/(X^{d_m+1}),\qquad
 \Gamma_m(\mathsf s)=
 \bigl(\mathsf s,E_{21/22}^{(d_m)}(\mathsf s;X)\bigr).
\]
The registered complete reader \(\Gamma_{E21}\) produces these representations.
For \(m'\ge m\), truncation
\(
 \tau_{m',m}(\mathsf s,[f]_{d_{m'}})
 =(\mathsf s,[f]_{d_m})
\)
satisfies
\begin{equation}
 \tau_{m',m}\circ\Gamma_{m'}=\Gamma_m,
 \qquad
 E_{21/22}:=\varprojlim_mE_{21/22}^{(d_m)}.
 \tag{A3}
 \label{eq:alpha-sistema-analitico-completo}
\end{equation}
No section of bare truncation is added. When the representation
of the next three coefficients is made explicit, it is the derived
projection
\[
 \Lambda_{d_m+1:d_m+3}
 :=\operatorname{pr}_{d_m+1:d_m+3}\circ
   \Gamma_{E21}\circ\mathcal U_{\mathrm{TPK}},
\]
not an independent generator or a reading of the target value.

For a cofinal set \(\mathcal N_\alpha\), the reader expresses a cofinal
function \(m:\mathcal N_\alpha\to\mathbb N\) and intervals
\(I_{B,N}=[\ell_N,u_N]\Subset I_\alpha\) such that
\begin{equation}
\begin{gathered}
 I_{B,N'}\subseteq I_{B,N}\ (N'\ge N),\qquad
 \operatorname{diam}I_{B,N}\longrightarrow0,\qquad
 I_{B,N}\subseteq C_N^A,\\
 E_{21/22}^{(d_{m(N)})}(\ell_N)\le0\le
 E_{21/22}^{(d_{m(N)})}(u_N),\qquad
 \inf_{x\in I_{B,N}}\partial_xE_{21/22}^{(d_{m(N)})}(x)>0,\\
 E_{21/22}^{(d_{m(N)})}\longrightarrow E_{21/22}
 \quad\text{uniformly on }I_\alpha .
\end{gathered}
 \tag{A4}
 \label{eq:alpha-via-b-proyectiva}
\end{equation}

\begin{proposition}[Compatibility squares and common cylinders]
\label{prop:alpha-cuadrados-compatibilidad-historico-rev04}
The restrictions of words and jets commute separately:
\[
\begin{array}{ccc}
 \mathcal X_\alpha^{\mathrm{enr}}&
 \xrightarrow{\ \mathsf T_{\alpha,N'}\ }&\mathcal W_{N'}\\[1mm]
 \Big\Vert&&\Big\downarrow\rho_{N',N}\\[1mm]
 \mathcal X_\alpha^{\mathrm{enr}}&
 \xrightarrow{\ \mathsf T_{\alpha,N}\ }&\mathcal W_N
\end{array}
\qquad
\begin{array}{ccc}
 \mathcal X_\alpha^{\mathrm{enr}}&
 \xrightarrow{\ \Gamma_{m'}\ }&\mathfrak J_{m'}\\[1mm]
 \Big\Vert&&\Big\downarrow\tau_{m',m}\\[1mm]
 \mathcal X_\alpha^{\mathrm{enr}}&
 \xrightarrow{\ \Gamma_m\ }&\mathfrak J_m .
\end{array}
 \tag{A5}
\]
The connection between the two squares is the certified inclusion
\(I_{B,N}\subseteq C_N^A\), with nesting, contraction, and cofinality.
The unique point
\(x_B:=\bigcap_{N\in\mathcal N_\alpha}I_{B,N}\) therefore satisfies
\begin{equation}
 \forall N\in\mathcal N_\alpha,\qquad
 \operatorname{pref}_N(x_B)=\alpha_{A,N}=:\alpha_{B,N}.
 \tag{A6}
 \label{eq:alpha-igualdad-cada-nivel}
\end{equation}
\end{proposition}

\begin{proof}
Compatibility of route A is
\eqref{eq:alpha-via-a-proyectiva-exacta}; that of route B is
\eqref{eq:alpha-sistema-analitico-completo}. Each jet has a simple root
in its interval by the change of sign and the derivative bound in
\eqref{eq:alpha-via-b-proyectiva}. The nested collection determines \(x_B\), and
uniform convergence gives \(E_{21/22}(x_B)=0\). Cofinal inclusion in the
cylinders of route A forces its prefixes. It has not been assumed that
truncating a polynomial preserves its root.
\end{proof}

\begin{theorem}[Projective and coinductive equality of the two routes]
\label{prop:alpha-dos-limites-historico-rev04}
For every compatible level of the common cofinal system,
\begingroup
\renewcommand{\theequation}{A7}
\begin{equation}
 \boxed{\alpha_{A,N}=\rho_N(\alpha_A)
 =\rho_N(\alpha_B)=\alpha_{B,N}.}
 \label{eq:alpha-dos-vias-cada-nivel-historico-rev04}
\end{equation}
\endgroup
In the coinductive limit,
\begingroup
\renewcommand{\theequation}{A8}
\begin{equation}
 \boxed{\alpha_A
 =\varprojlim_N\alpha_{A,N}
 =\varprojlim_N\alpha_{B,N}
 =\operatorname*{arg\,zero}_{x\in I_\alpha}E_{21/22}(x)
 =\alpha_B=\alpha_{\mathrm{HMT}}.}
 \label{eq:alpha-dos-vias-limite-coinductivo-historico-rev04}
\end{equation}
\endgroup
No conventional value enters: metrological recognition follows
the two internal representations.
\end{theorem}

\begin{proof}
Route A forms compatible cylinders of diameter \(1000^{-N}\); route B,
nested intervals of diameter tending to zero. By
\eqref{eq:alpha-via-b-proyectiva}, the point of the second collection is the
simple and unique root of the complete kernel. Since it belongs to a cofinal
subfamily of the cylinders of the first, both intersections contain the same point.
Equation~\eqref{eq:alpha-igualdad-cada-nivel} gives equality at every
compatible depth, and the inverse limit gives (A8). Finite executions
certify only projections; they do not bound the statement.
\end{proof}

The simple root \(\xi_9\) and the witness
\eqref{eq:alpha-testigo-finito-nueve} are finite verifications of this
system. They neither define the limit nor reduce (A7)--(A8) to a truncated
coincidence: complete equality follows from the state extensions, the two
compatibility squares, and uniqueness of the cofinal intersection.
